\section{Simplicity, persistence, and one fixed metric}
\label{sec:spectral}

The preceding construction produces an exact eigenfunction with a large
sign score.  We first make its eigenvalue simple while preserving that
function exactly.  The freedom required for this step is supplied by the
three directions tangent to its level sets in an untouched part of $P$.

\begin{lemma}[Simplicity preserving an eigenfunction]
\label{spec:simple}
Let $g$ be a smooth metric on $M$, and suppose
\[
 -\Delta_g t=\Lambda_k t,\qquad k\ge1,
\]
\[
 g=g_0,\quad t=u_k^0
 \quad\hbox{outside a compact subset of }P.
\]
In every smooth-topology neighborhood of $g$ there is a smooth metric
$\widetilde g$ for which $\Lambda_k$ is simple, with the same
eigenfunction $t$.  Moreover, $\widetilde g-g$ has compact support in
$P$, and $d\operatorname{vol}_{\widetilde g}=d\operatorname{vol}_g$.
\end{lemma}

\begin{proof}
We shall use only perturbations in open subsets of $P$ where the metric
and eigenfunction still agree with the background.  Such a region remains
after any finite number of compactly supported changes.

First, the squared gradient norm of the background varies along its level sets.
More precisely, on every nonempty open subset of $P$ there is a smaller
open subset on which $dt$ and $d|dt|_{g_0}^2$ are independent when
$t=u_k^0$.  Indeed, writing $\phi=\theta+z$, Equation~\eqref{env:background-gradient-identity} gives
\[
 |du_k^0|_{g_0}^2
 =k^2\bigl(H_k(p)-2(u_k^0)^2\bigr),\qquad
 H_k(p)=\sin^{2k}p+\sin^{2k-2}p.
\]
The first summand in $H_k$ has strictly positive derivative on
$0<p<\pi/2$; the second is increasing for $k>1$ and constant for
$k=1$.  Thus $H_k'(p)>0$ for every $k\ge1$.  Wherever
$\sin(k\phi)\ne0$, the covectors $du_k^0$ and $dp$ are independent, so
\[
 d|du_k^0|_{g_0}^2\wedge du_k^0
 =k^2 H_k'(p)\,dp\wedge du_k^0\ne0.
\]
The excluded set has empty interior.

Let $E=\ker(-\Delta_g-\Lambda_k)$ have dimension $m>1$, and choose a
real $t_2\in E$ independent of $t$.  Take an open background region
$V\Subset P$ on which $dt$ and $d|dt|_g^2$ are independent.  We claim
that the tangential component of $\grad_g t_2$ along a level set of $t$
is nonzero somewhere in $V$.  Otherwise $dt_2$ vanishes on $\ker dt$
throughout $V$.  Since $dt\ne0$, submersion coordinates on a smaller
product neighborhood give $t_2=F(t)$ for a smooth one-variable function
$F$.  The two eigenfunction equations imply
\begin{equation}
 F''(t)|dt|_g^2=\Lambda_k\bigl(tF'(t)-F(t)\bigr).
 \label{spec:local-dependence}
\end{equation}
On a still smaller neighborhood there is a smooth vector field $X$ with
$Xt=0$ and $X(|dt|_g^2)\ne0$.  Differentiating
\eqref{spec:local-dependence} along $X$ gives $F''(t)=0$ there.
The image of this neighborhood under the submersion $t$ contains an
open interval; hence $F(r)=ar+b$ on that interval.  Substitution into
\eqref{spec:local-dependence}, using $\Lambda_k>0$, forces $b=0$.
Consequently $t_2-at$ vanishes on a nonempty open set.  Weak unique
continuation for smooth scalar elliptic equations on a connected
manifold, a consequence of Aronszajn's Theorem
\cite[pp.~235--236]{Aronszajn1957}, makes it identically zero.
The hypotheses apply because the metric is smooth and positive and the
zeroth-order coefficient $\Lambda_k$ is smooth.  This contradicts the
choice of $t_2$ and proves the claim.

On a small neighborhood of a point supplied by the claim, let $\Pi$ be
the $g$-orthogonal projection onto $\ker dt$ and put
$z=\Pi\grad_g t_2$.  Choose a nonzero nonnegative smooth bump $\chi$
supported in that neighborhood, with $z\ne0$ on its support.  Extend by
zero the endomorphism
\begin{equation}
 J=\chi\left(z\otimes z^{\flat_g}-\frac{|z|_g^2}{3}\Pi\right).
 \label{spec:tangential-variation}
\end{equation}
It is smooth, $g$-self-adjoint, trace free, and annihilates $\grad_g t$.
Here the factor $1/3$ uses $\operatorname{rank}\Pi=3$.  Moreover,
\begin{equation}
 \mathfrak b_J(t_2,t_2)
 :=\int_M g(J\grad_g t_2,\grad_g t_2)\,
                  d\operatorname{vol}_g
 =\frac23\int_M\chi|z|_g^4\,d\operatorname{vol}_g>0.
 \label{spec:positive-variation}
\end{equation}
For real $s$ define an ordinary metric by
\begin{equation}
 g_s^{-1}=\exp(sJ)g^{-1}.
 \label{spec:metric-path}
\end{equation}
Self-adjointness makes this a positive symmetric cometric.  Since
$\det\exp(sJ)=\exp(s\operatorname{tr}J)=1$, its volume density equals
that of $g$.  Since $J\grad_g t=0$, it also satisfies
$g_s^{-1}dt=g^{-1}dt$.  The divergence formula for the Laplacian
therefore gives
\[
 -\Delta_{g_s}t=\Lambda_k t
 \quad\hbox{for every }s.
\]
Thus both volume and the flux of $dt$ are preserved exactly.

We show that the multiplicity at $\Lambda_k$ is less than $m$ for all
sufficiently small nonzero $s$.  Suppose instead that $s_n\to0$,
$s_n\ne0$, and that there are $m$ real orthonormal eigenfunctions
$e_{1,n},\ldots,e_{m,n}$ of $g_{s_n}$ at $\Lambda_k$.  All norms and
orthogonality may be taken with the same volume density.  Smooth
convergence of the metrics, elliptic estimates, and a diagonal compactness
argument give a subsequence such that $e_{a,n}\to e_a$ smoothly for
$1\le a\le m$.  To spell out the uniform estimate used here, for every
fixed integer $r\ge0$, comparison with the limiting operator and
absorption of its small coefficient difference give
\[
 \|v\|_{H^{r+2}}
 \le C_r\bigl(\|\Delta_{g_{s_n}}v\|_{H^r}
                         +\|v\|_{L^2}\bigr)
\]
for all sufficiently large $n$.  Iteration at the fixed eigenvalue gives
uniform Sobolev bounds of every order; compact Sobolev embeddings give
the asserted convergence.  The limits form an orthonormal basis of $E$.

For $v\in E$, the weak equation for $e_{a,n}$ with metric $g_{s_n}$
and the weak equation for $v$ with metric $g$ have the same right-hand
side.  Subtracting and dividing by $s_n$ yields
\[
 \int_M g\left(\frac{\exp(s_nJ)-I}{s_n}\grad_g e_{a,n},
                    \grad_g v\right)\,d\operatorname{vol}_g=0.
\]
Taking the limit shows $\mathfrak b_J(e_a,v)=0$ for all $a$ and all
$v\in E$.  Since the $e_a$ span $E$, this contradicts
\eqref{spec:positive-variation}.  This proves the strict decrease.

The eigenfunction $t$ survives every decrease, so the multiplicity is
always at least one.  At most $m-1$ repetitions make it one.  At each
repetition the parameter can be taken arbitrarily small in the smooth
topology.  For example, choose a finite $C^r$ tolerance contained in the
prescribed neighborhood and allocate to these finitely many changes a
total tolerance smaller than it.  Their supports form a finite union of
compact subsets of $P$, leaving another background region whenever one
is needed.  Volume and $t$ are preserved throughout.  This proves all
the assertions.
\end{proof}

Fix a constant $c_{\mathrm H}>0$ such that the comparison in
Lemma~\ref{score:measure} reads
\begin{equation}
 \mathcal H_g^3(\{v=0\})\ge c_{\mathrm H}\mathcal S(v),\qquad g\in\mathcal U,
 \label{spec:score-comparison}
\end{equation}
for every continuous real $v$ and every sign score used in the construction.
Its uniformity follows from the fixed metric comparability on
$\mathcal U$.  When considering persistence, the cubes and displacements
defining $\mathcal S$ are held fixed.

\begin{lemma}[Open persistence of a fixed score]
\label{spec:persistence}
Suppose $g\in\mathcal U$ has a simple positive eigenvalue $\lambda$
with real eigenfunction $t$, and, for a fixed sign score and $N>0$,
\[
 \lambda>N,\qquad c_{\mathrm H}\mathcal S(t)>N\sqrt\lambda.
\]
There is a smooth-topology neighborhood $\mathcal V\subset\mathcal U$
of $g$ such that every $g'\in\mathcal V$ has a simple eigenvalue
$\lambda'>N$ with a real eigenfunction $t'$ satisfying
\[
 c_{\mathrm H}\mathcal S(t')>N\sqrt{\lambda'}.
\]
In particular its nodal-measure ratio exceeds $N$.
\end{lemma}

\begin{proof}
Choose a number $T$ and an open interval $I$ containing $\lambda$ such that
\[
 N\sqrt\lambda<T<c_{\mathrm H}\mathcal S(t),
 \qquad I\subset\bigl(N,(T/N)^2\bigr).
\]
These choices are possible by the two strict inequalities in the
hypothesis. Normalize $t$ in $L^2(g)$; scalar invariance leaves its
score unchanged. The functional
$\mathcal F=c_{\mathrm H}\mathcal S$ is lower semicontinuous for
uniform convergence and invariant under nonzero real rescaling by
\eqref{score:lsc} and the paragraph following it.
Apply Lemma~\ref{lem:simple-eigenpair-continuity} to this functional,
the threshold $T$, and the interval $I$, and intersect its neighborhood
with $\mathcal U$. For every metric $g'$ in the resulting neighborhood,
the selected simple eigenpair satisfies
\[
 \lambda'\in I,\qquad c_{\mathrm H}\mathcal S(t')>T
                         >N\sqrt{\lambda'}.
\]
In particular $\lambda'>N$, and
\eqref{spec:score-comparison} gives the nodal-measure conclusion.
\end{proof}

\begin{proposition}[An arbitrarily close step]
\label{spec:one-step}
Let $h\in\mathcal U$ and suppose $h-g_0$ has compact support in $P$.
For every smooth-topology neighborhood $\mathcal O$ of $h$, every
$N>0$, and every integer $k_{\min}$, there are an integer $k\ge k_{\min}$, a metric
$\widetilde h\in\mathcal O\cap\mathcal U$, and a real eigenfunction
$t$ such that
\[
 -\Delta_{\widetilde h}t=\Lambda_k t,\qquad
 \Lambda_k>N,\qquad c_{\mathrm H}\mathcal S(t)>N\sqrt{\Lambda_k}
\]
for a sign score $\mathcal S$, and $\Lambda_k$ is simple.
Both $\widetilde h-g_0$ and $t-u_k^0$ have compact support in $P$.
The inequalities persist in a neighborhood of $\widetilde h$ in the
sense of Lemma~\ref{spec:persistence}.
\end{proposition}

\begin{proof}
First choose a finite $C^r$ neighborhood of $h$ contained in
$\mathcal O\cap\mathcal U$, reserving part of its tolerance for
Lemma~\ref{spec:simple}.  Let $c_{\mathrm s}>0$ be the constant,
independent of the amplified envelope, in the conclusion of
Proposition~\ref{score:good}.  By Lemma~\ref{env:gaps}, choose the nested
patches and an envelope with the required gaps and with
\[
 c_{\mathrm H}c_{\mathrm s}\int_D B(x)\,dx>2\sqrt3\,N.
\]
For these fixed data, Lemma~\ref{corr:realization} specifies finite residual derivative orders and
accuracy powers that suffice for the reserved $C^r$ tolerance.
Construct the packets to those orders and use
Proposition~\ref{score:good}, for arbitrarily large integers $k$, to obtain
a real $u$ satisfying the residual and noncancellation hypotheses as well
as
\[
 \mathcal S(u)\ge c_{\mathrm s}k\int_D B(x)\,dx.
\]
Choose $k\ge k_{\min}$ with $\Lambda_k>N$ above all these thresholds and those
required by the correction.  All order choices precede this choice of $k$.
Lemma~\ref{corr:realization} gives an ordinary metric $g_*$ within the reserved
tolerance and an exact $\Lambda_k$ eigenfunction $t$, a positive smooth
multiple of $u$, with the stated background identities outside a compact
subset of $P$.  Positive-factor invariance gives
$\mathcal S(t)=\mathcal S(u)$.  Since
$\sqrt{\Lambda_k}=\sqrt{2k^2+k}\le\sqrt3\,k$ for $k\ge1$,
\[
 c_{\mathrm H}\mathcal S(t)>2N\sqrt{\Lambda_k}.
\]
Lemma~\ref{spec:simple} then makes
$\Lambda_k$ simple by an arbitrarily small change preserving $t$.
Take this change within the remaining tolerance.  All support conditions
are preserved, and Lemma~\ref{spec:persistence} supplies the last claim.
\end{proof}

\begin{proof}[Proof of Theorem~\ref{thm:four}]
Let $\mathcal X=C^\infty(M;\operatorname{Sym}^2T^*M)$, the real
Fr\'echet space of smooth symmetric covariant two-tensors.  Fix $C^r$
seminorms defined using $g_0$ and the complete translation-invariant metric
\[
 d_{\mathcal X}(a,b)
 =\sum_{r=0}^{\infty}2^{-r-1}
                 \min\{1,\|a-b\|_{C^r}\}.
\]
It induces the smooth topology.  Completeness follows because a Cauchy
sequence converges in every $C^r$ norm and the limits agree under the
natural inclusions.  The positive metrics form an open subset of
$\mathcal X$, and $\mathcal U$ is open in this ambient space.

Take any sequence $N_j\to\infty$, with $N_j>0$.  Starting from $g_0$,
apply Proposition~\ref{spec:one-step} with threshold $N_1$ to choose
a first center $h_1$ and
its persistence neighborhood $\mathcal V_1\subset\mathcal U$.
Choose a closed $d_{\mathcal X}$-ball $\mathcal B_1$ of positive radius
centered at $h_1$, contained in $\mathcal V_1$, and of diameter at most
$2^{-1}$.  Inductively, suppose $h_{j-1}$ and $\mathcal B_{j-1}$ have
been chosen.  Its interior is a neighborhood of $h_{j-1}$.
Proposition~\ref{spec:one-step}, with threshold $N_j$, therefore gives
a center
\[
 h_j\in\operatorname{int}\mathcal B_{j-1}
\]
whose difference from $g_0$ has compact support in $P$, together with a
persistence neighborhood $\mathcal V_j$.  Choose a positive-radius
closed ball centered at $h_j$ such that
\[
 \mathcal B_j\subset
 \operatorname{int}\mathcal B_{j-1}\cap\mathcal V_j\cap\mathcal U,
 \qquad \operatorname{diam}\mathcal B_j\le2^{-j}.
\]
This is possible because the set on the right is open and contains
$h_j$.  Every open neighborhood in the smooth topology contains a
neighborhood specified by finitely many seminorms, exactly the closeness
provided by Proposition~\ref{spec:one-step}.

These are nested nonempty closed sets with diameters tending to zero in
the complete space $\mathcal X$.  Their centers are Cauchy, and their
limit $g_\infty$ belongs to every $\mathcal B_j$.  In particular
$g_\infty\in\mathcal U$,
so it is a smooth positive Riemannian metric on the fixed closed connected
four-manifold $M=S^2\times\mathbb T^2$.  For each $j$, membership in
$\mathcal V_j$ supplies an exact nonzero real eigenfunction $u_j$ of this
same metric, with eigenvalue $\lambda_j$ satisfying
\[
 \lambda_j>N_j,\qquad
 \frac{\mathcal H_{g_\infty}^3(\{u_j=0\})}{\sqrt{\lambda_j}}
 \ge\frac{c_{\mathrm H}\mathcal S_j(u_j)}{\sqrt{\lambda_j}}>N_j.
\]
Here $\mathcal S_j$ is the fixed score from the $j$th step.
Both quantities therefore tend to infinity.  Compact support in $P$
was needed only for the finite-stage centers; it is not a hypothesis on
the limit used in this conclusion.  Likewise the final eigenvalues may
move from the intermediate values $\Lambda_k$: persistence gives the
exact eigenpairs and the two strict inequalities at the single limiting
metric.  This proves the stated refutation.
\end{proof}
